% ============================================
% Open-Rank Faculty Cover Letter
% Florida Polytechnic University
% ============================================

\documentclass[11pt,letterpaper]{letter}

\usepackage{geometry}
\usepackage{microtype}
\usepackage[hidelinks]{hyperref}

\geometry{left=25mm,right=25mm,top=25mm,bottom=25mm}

\signature{Decory Edwards}
\address{Department of Economics \\ Johns Hopkins University \\ Baltimore, MD 21218 \\ \texttt{dedwar65@jh.edu}}

\begin{document}

\begin{letter}{Search Committee \\
Florida Polytechnic University}

\opening{Dear Members of the Search Committee,}

I am writing to apply for the open-rank faculty position in Business Analytics, Engineering Management, and Data Science at Florida Polytechnic University. I am a Ph.D. candidate in Economics at Johns Hopkins University, advised by Professors Christopher D. Carroll, Nicholas W. Papageorge, and Stelios Fourakis, and I expect to complete my degree in 2026. My research and teaching interests lie in macroeconomics, inequality, and applied quantitative methods, with a strong emphasis on data-driven analysis and policy-relevant questions.

My research agenda focuses on understanding wealth inequality through the lens of heterogeneous returns to wealth and household behavior. In my job market paper, I use microdata from the Survey of Consumer Finances to estimate persistent heterogeneity in portfolio returns and embed these estimates in a structural macroeconomic model. I show that allowing for return heterogeneity substantially improves the model’s ability to match observed wealth concentration and generates realistic implications for consumption behavior and policy analysis. This work connects micro-level financial behavior to aggregate outcomes in a way that is empirically grounded and directly relevant for applied analysis.

In related work, I use data from the Health and Retirement Study to study how trust and beliefs shape household investment outcomes. I find evidence of a nonlinear relationship between trust and realized returns, highlighting a behavioral channel through which attitudes and expectations influence economic outcomes. Together, my research emphasizes careful data construction, transparent empirical methods, and the interpretation of complex quantitative results.

\textbf{Teaching.} My teaching experience centers on undergraduate macroeconomics and applied quantitative courses. I have led weekly discussion sections and held regular office hours for Principles of Macroeconomics, Intermediate Macroeconomic Theory, and an upper-level macroeconomic policy seminar. My teaching emphasizes clarity, applied problem-solving, and the use of real data to motivate theory. I aim to equip students with the ability to work with data, interpret empirical results, and communicate findings clearly. I am prepared to teach courses that emphasize data analysis, applied economics, and quantitative decision-making.

Florida Polytechnic University’s mission as a STEM-focused institution with a strong emphasis on applied learning makes it a particularly good fit for my background. The Business Analytics and Data Science programs emphasize the use of quantitative tools to solve real-world problems, an approach that closely aligns with my training in empirical modeling, simulation methods, and large-scale data analysis. My research and teaching experience would allow me to contribute to courses that integrate economic reasoning with data analytics, such as applied modeling, predictive analysis, and the evaluation of economic and business outcomes using real datasets. I am especially drawn to Florida Poly’s emphasis on interdisciplinary collaboration and on preparing students for careers where analytical rigor and practical relevance are equally important.

I would welcome the opportunity to contribute to Florida Polytechnic University’s teaching mission, curriculum development, and applied research initiatives. Thank you for your time and consideration. I look forward to the possibility of discussing my application with you.

\closing{Sincerely,}

\end{letter}
\end{document}
